% !TeX root = ../../Book3.tex
% 纲要标题：## 第23章*　Hensel 性、非分歧与平展
% 纲要位置：计划纲要.md，第 2218 行。

\OptionalChapter{Hensel 性、非分歧与平展}{b3:ch:23}
\BookChapterNumber{b3:ch:23}

\begin{chapterabstract}
本章从约化方程的单根和互素因式出发，证明完备局部环中的 Hensel 提升，并说明 Hensel 化与完备化所解决的问题不同。随后以平方零理想上的提升为依据，建立形式光滑的关系模分裂判据、光滑映射的局部标准形以及平展的等价刻画。最后逐项计算尖点正规化的微分模、纤维长度和原点外的逆映射，区分正则性、光滑性、非分歧性与平坦性。
\end{chapterabstract}

\begin{learninggoals}
\begin{itemize}
\item 用误差修正证明单根提升与互素因式提升，并辨认导数和互素条件的作用。
\item 根据泛性质区分 Hensel 化与完备化，理解相同的有限阶商不意味着环相同。
\item 由平方零提升建立关系模的分裂，并把有限呈示转化为邻域上的有限方程组。
\item 证明平展的几种刻画等价，区分平坦性与微分模消失这两个条件。
\item 完成尖点正规化的三个局部计算，说明光滑性为何必须连同底环一起讨论。
\end{itemize}
\end{learninggoals}

本章为选读。\chapref{b3:ch:13}的完备化、\chapref{b3:ch:14}的微分与形式提升，以及\chapref{b3:ch:12}的平坦性是主要准备；局部性质比较还使用\chapref{b3:ch:18}的正则性。第14章已经给出形式光滑、形式非分歧、形式平展和有限呈示条件下的定义，本章集中证明它们之间尚未建立的联系，不重新组织这些基础定义。

\ProseParagraphBreak
单根与因式提升、关系模判据、局部标准形、平展的等价刻画及尖点计算均在正文证明。三项较深结果明确作为外部引用：Hensel 化的存在及基本性质、Noether 局部形式光滑蕴含平坦、完美基域上的正则与光滑比较。第一项只用于\secref{b3:sec:2301}，第二项用于一般标准模型的平坦性，第三项只用于最后一节的比较；它们不互相充当证明。

\ProseParagraphBreak
Hensel 因式提升见松村英之\cite[Theorem 8.3, p. 58]{Matsumura1989CommutativeRingTheory}，Hensel 化的构造与泛性质见永田雅宜\cite[\S43, (43.3), (43.5), (43.8)--(43.10), pp. 180--183]{Nagata1975LocalRings}。形式光滑与关系模的联系可读松村英之\cite[Theorem 25.2, pp. 194--195]{Matsumura1989CommutativeRingTheory}；本章借用的局部平坦性定理及正则性比较，分别定位于同书\cite[Theorem 28.9, p. 222; Theorem 30.3 and Remark 2, pp. 233--235]{Matsumura1989CommutativeRingTheory}。

% !TeX root = ../../Book3.tex
% 纲要标题：### 23.1 Hensel 引理
% 纲要位置：计划纲要.md，第 2220 行。

\section{Hensel 引理}
\label{b3:sec:2301}

% 纲要内容：
%
% * 单根提升与因式分解提升
% * 完备局部环的 Hensel 性
% * Hensel 化和完备化的区别
%

\subsection{单根提升与因式分解提升}
\label{b3:subsec:2301:01}

设 \((R,\mathfrak m)\) 是局部环，\(k=R/\mathfrak m\)。把多项式的系数约化到 \(k\)，常会得到一个容易求解的方程。问题是：约化方程的解，能否成为原方程的解的第一项？如果约化根是单根，导数便可以用来修正误差；如果它是重根，误差未必能消去，也可能有不止一种消去方法。

\begin{theorem}[单根的 Hensel 提升]\label{b3:thm:hensel-simple-root}
设 \((R,\mathfrak m)\) 为交换含幺局部环，且自然映射 \(R\to\varprojlim_n R/\mathfrak m^n\) 是同构。设 \(f\in R[T]\)，\(\bar a\in k=R/\mathfrak m\) 满足
\[
\bar f(\bar a)=0,\qquad \bar f'(\bar a)\ne0.
\]
则存在唯一的 \(a\in R\)，使 \(f(a)=0\) 且 \(a\bmod\mathfrak m=\bar a\)。这里不要求 \(f\) 首一。
\end{theorem}
\begin{proof}
任取 \(\bar a\) 的提升 \(a_0\)。因为 \(f'(a_0)\) 模 \(\mathfrak m\) 非零，它是单位。递归令
\[
a_{j+1}=a_j-\frac{f(a_j)}{f'(a_j)}.
\]
每次修正都在 \(\mathfrak m\) 中，故所有 \(a_j\) 具有相同的剩余类，\(f'(a_j)\) 始终可逆。多项式恒等式
\(f(u+v)=f(u)+f'(u)v+v^2h(u,v)\) 给出：若 \(f(a_j)\in\mathfrak m^{2^j}\)，则
\[
a_{j+1}-a_j\in\mathfrak m^{2^j},\qquad
f(a_{j+1})\in\mathfrak m^{2^{j+1}}.
\]
初始误差属于 \(\mathfrak m\)，所以归纳成立。由完备性，\((a_j)\) 收敛到 \(a\in R\)。有限次加法与乘法对 \(\mathfrak m\) 进拓扑连续，故 \(f(a)\in\bigcap_n\mathfrak m^n=0\)，并且 \(a\) 仍提升 \(\bar a\)。

若 \(a,b\) 都是这样的根，则
\[
0=f(a)-f(b)=(a-b)\bigl(f'(b)+(a-b)c\bigr)
\]
其中 \(c\in R\)。括号内的元素模 \(\mathfrak m\) 等于 \(\bar f'(\bar a)\)，所以是单位，因而 \(a=b\)。这个唯一性论证本身不需要完备性。
\end{proof}

\begin{example}\label{b3:exm:hensel-root-five-adic}
在 \(\mathbb Z_5\) 中，\(T^2+1\) 的约化根 \(2\) 是单根。故恰有一个根同余于 \(2\pmod5\)。写作 \(2+5c\)，模 \(25\) 的方程为 \(5+20c\equiv0\)，得 \(c\equiv1\pmod5\)，所以这个根模 \(25\) 等于 \(7\)。从约化根 \(3\) 得到另一个根，它是前一个的负数。这里是在指定剩余类内唯一，并不是说多项式在环中只能有一个根。
\end{example}

\begin{definition}\label{b3:def:henselian-local-ring}
设 \((R,\mathfrak m,k)\) 为交换含幺局部环。若对每个首一 \(f\in R[T]\) 及每个分解
\[
\bar f=\bar g\bar h\quad\text{于 }k[T],
\]
其中 \(\bar g,\bar h\) 首一且互素，都存在首一 \(g,h\in R[T]\)，满足
\(f=gh\)、\(g\bmod\mathfrak m=\bar g\)、\(h\bmod\mathfrak m=\bar h\)，并且次数分别等于 \(\deg\bar g,\deg\bar h\)，则称 \(R\) 为 \term[Hensel-huan]{Hensel 环}[Henselian local ring]，也称 \(R\) 具有 Hensel 性。
\end{definition}

在这个定义中取 \(\bar g=T-\bar a\)，便得到首一多项式的单根提升：另一个因子与它互素，恰好表示 \(\bar f'(\bar a)\ne0\)。因式提升还可以保留一整组根，而不要求它们分别属于剩余域。下面直接证明完备性足以保证这种较一般的提升。

\subsection{完备局部环的 Hensel 性}
\label{b3:subsec:2301:02}

\begin{theorem}[互素因式的 Hensel 提升]\label{b3:thm:hensel-factor-lifting}
设 \((R,\mathfrak m,k)\) 为交换含幺局部环，并且 \(R\cong\varprojlim_n R/\mathfrak m^n\)。设 \(f\in R[T]\) 首一，且
\(\bar f=\bar g\bar h\)，其中 \(\bar g,\bar h\in k[T]\) 首一互素，次数分别为 \(r,s\)。则存在唯一的首一分解 \(f=gh\)，其中 \(g,h\) 具有上述次数和约化像。特别地，每个完备分离局部环都是 Hensel 环。
\end{theorem}
\begin{proof}
可假设 \(r,s>0\)，因为零次的首一因子只能是 \(1\)。先在 \(k\) 上考虑线性映射
\[
L:k[T]_{<r}\oplus k[T]_{<s}\longrightarrow k[T]_{<r+s},
\qquad (u,v)\longmapsto\bar h u+\bar g v.
\]
若像为零，由 \(\bar g\mid\bar h u\) 及互素性可得 \(\bar g\mid u\)。次数限制迫使 \(u=0\)，继而 \(v=0\)。两端维数均为 \(r+s\)，故 \(L\) 为同构。这正是每一阶修正都可以且只能有一种方式完成的原因。

任取首一提升 \(g_1,h_1\)，次数固定为 \(r,s\)，于是 \(f-g_1h_1\in\mathfrak m R[T]\)。设已构造 \(g_n,h_n\)，使误差属于 \(\mathfrak m^nR[T]\)。在 \(\mathfrak m^n/\mathfrak m^{n+1}\) 上对 \(L\) 作张量积，仍为同构，故存在系数属于 \(\mathfrak m^n\) 的多项式 \(u_n,v_n\)，次数分别小于 \(r,s\)，使
\[
f-g_nh_n\equiv h_nu_n+g_nv_n\pmod{\mathfrak m^{n+1}R[T]}.
\]
令 \(g_{n+1}=g_n+u_n\)、\(h_{n+1}=h_n+v_n\)。乘积中另外出现的 \(u_nv_n\) 属于 \(\mathfrak m^{2n}R[T]\subseteq\mathfrak m^{n+1}R[T]\)，故新的误差提高一阶。修正不改变首项和次数。

固定次数使每个系数各自成为 Cauchy 序列。逐个取极限得到首一 \(g,h\)；分离性使 \(f-gh=0\)。若还有另一个分解 \(g'h'\)，起初 \(g'-g,h'-h\) 的系数在 \(\mathfrak m\) 中。假设已经在 \(\mathfrak m^n\) 中，则模 \(\mathfrak m^{n+1}\) 有
\[
\bar h(g'-g)+\bar g(h'-h)=0.
\]
由同一个 \(L\) 的单射性，两差的系数都在 \(\mathfrak m^{n+1}\) 中。归纳并使用分离性可得 \(g'=g,h'=h\)。
\end{proof}

上述证明采用逐次修正系数的方法；同样的因式提升可参阅松村英之的《Commutative Ring Theory》\cite[Theorem 8.3, p. 58]{Matsumura1989CommutativeRingTheory}。单根情形的 Newton 修正每次至少使误差的理想次幂加倍，而这里逐阶求解一个固定的线性方程组；两者都依靠一次近似处的可逆性。

\ProseParagraphBreak
互素条件不能省略。例如 \(T^2-5\) 在 \(\mathbb Z_5\) 中约化为 \(T\cdot T\)，但若有根 \(a\)，则 \(2v_5(a)=1\)，不可能成立。因此这个约化分解不能提升为两个首一一次因子。完备性保证误差修正的极限存在，却不能代替求解修正方程所需的可逆条件。

\subsection{Hensel 化与完备化}
\label{b3:subsec:2301:03}

完备化同时加入所有相容的 \(\mathfrak m\) 进近似。若只希望使上述互素因式能够提升，就不必加入这么多元素。描述这个较小扩张的合适方式，是规定它到其他 Hensel 局部环的映射。

\begin{definition}\label{b3:def:henselization}
设 \((R,\mathfrak m)\) 为交换含幺局部环，\(\iota:R\to R^{\mathrm h}\) 为到 Hensel 局部环的局部同态。若对任意 Hensel 局部环 \(S\) 和局部同态 \(u:R\to S\)，存在唯一的局部同态 \(v:R^{\mathrm h}\to S\)，使 \(v\iota=u\)，则称 \(\iota:R\to R^{\mathrm h}\) 为 \(R\) 的 \term[Hensel-hua]{Hensel 化}[Henselization]。
\end{definition}

这一定义立即给出唯一性：两个具有该性质的对象之间各有唯一的相容局部同态，复合由唯一性等于恒等映射。存在性则不是形式推论；它需要把带有指定剩余点的有限代数扩张组织起来。下面明确引用这一构造定理，不把完备化当作它的证明。

\begin{theorem}[Hensel 化的存在；外部引用]\label{b3:thm:henselization-existence}
每个交换含幺局部环 \((R,\mathfrak m,k)\) 都有 Hensel 化 \(R^{\mathrm h}\)。它的极大理想为 \(\mathfrak mR^{\mathrm h}\)，剩余域为 \(k\)，且 \(R^{\mathrm h}\) 在 \(R\) 上忠实平坦。若 \(R\) 为 Noether 环，则 \(R^{\mathrm h}\) 也为 Noether 环，而且对每个 \(n\ge1\)，自然映射
\[
R/\mathfrak m^n\longrightarrow
R^{\mathrm h}/\mathfrak m^nR^{\mathrm h}
\]
为同构。因此二者的相应完备化自然同构。
\end{theorem}

本定理的存在性、泛性质、平坦性及完备化比较，采用永田雅宜《Local Rings》中的结论\cite[\S43, (43.3), (43.5), (43.8)--(43.10), pp. 180--183]{Nagata1975LocalRings}。该书先在正规情形构造扩张，再由商环处理一般情形；其中局部平坦同态的忠实性也可由本卷忠实平坦判据得出。本节不展开整个存在性证明。它是本选读章所引用的独立结果，后面光滑与平展的证明不以此定理为前提。

\ProseParagraphBreak
对 Noether 局部环，由 \(\widehat R\) 的 Hensel 性及泛性质，得到相容的局部同态
\[
R\longrightarrow R^{\mathrm h}\longrightarrow\widehat R.
\]
前者只需解决 Hensel 提升问题，后者还允许任意阶的相容极限。不能仅凭两环有相同的各阶商，就断言两环相同；还必须检查它们自身是否完备。

\begin{example}\label{b3:exm:henselization-not-completion}
令 \(R=\mathbb Q[t]_{(t)}\)，则 \(\widehat R=\mathbb Q[[t]]\)。多项式 \(F(X)=X^2-X-t\) 模 \(t\) 有单根 \(0\)，但在 \(\mathbb Q(t)\) 中没有根：若有根，则 \(1+4t\) 是有理函数的平方；然而它在 \(t=-1/4\) 处的零点阶数为一，平方的零点阶数必为偶数。因此 \(R\) 不是 Hensel 环。

\ProseParagraphBreak
在 Hensel 化中，这个根存在；它在完备化中的展开以
\[
-t+t^2-2t^3+5t^4+\cdots
\]
开始。这里 \(R^{\mathrm h}\) 可以看作 \(\widehat R\) 的子环：前者为 Noether 局部环，它到自身完备化的映射由 Krull 交定理为单射，而这个完备化正是 \(\widehat R\)。进一步，由永田书中 (43.9)，Hensel 化的每个元素都代数于 \(\mathbb Q(t)\)，所以 \(R^{\mathrm h}\) 是可数集；\(\mathbb Q[[t]]\) 却不可数，因为它包含所有系数只取 \(0,1\) 的无穷序列。因此
\[
R\subsetneq R^{\mathrm h}\subsetneq\widehat R.
\]
这个例子同时区分了局部化、Hensel 化与完备化，而不是仅仅赋予同一环三个名称。
\end{example}

% !TeX root = ../../Book3.tex
% 纲要标题：### 23.2 形式性质
% 纲要位置：计划纲要.md，第 2226 行。

\section{形式性质}
\label{b3:sec:2302}

% TODO: 按以下纲要要点撰写本节。

% 纲要内容（仅作写作计划，尚非教材正文）：
%
% * 平方零扩张上的提升问题
% * 形式光滑、形式非分歧和形式平展
% * 有限呈示条件的作用
%

待撰写。

% !TeX root = ../../Book3.tex
% 纲要标题：### 23.3 光滑与平展
% 纲要位置：计划纲要.md，第 2232 行。

\section{光滑与平展}
\label{b3:sec:2303}

% TODO: 按以下纲要要点撰写本节。

% 纲要内容（仅作写作计划，尚非教材正文）：
%
% * 光滑映射的标准局部模型
% * 平展的平坦性及非分歧性
% * 不把局部同胚直觉替代代数条件
%

\subsection{光滑映射的标准局部模型}
\label{b3:subsec:2303:01}

待撰写。

\subsection{平展映射的平坦性与非分歧性}
\label{b3:subsec:2303:02}

待撰写。

\subsection{局部同胚直觉与平展的代数条件}
\label{b3:subsec:2303:03}

待撰写。

% !TeX root = ../../Book3.tex
% 纲要标题：### 23.4 局部性质的比较与反例
% 纲要位置：计划纲要.md，第 2238 行。

\section{局部性质的比较与反例}
\label{b3:sec:2304}

% 纲要内容：
%
% * 尖点正规化在尖点处的局部计算
% * 光滑性、正则性与可分性的比较
% * 代数几何中几种拓扑的简介
% * 在特征零下令 \(A=k[t^2,t^3]\subseteq B=k[t]\)，计算 \(\Omega_{B/A}\cong B/(t)\,dt\)，故正规化在尖点处并非非分歧
% * 计算尖点纤维 \(B/(t^2,t^3)B\cong k[t]/(t^2)\)；其长度为2而双有理映射的泛秩为1，用有限平坦模在局部环上自由的结论证明尖点处不平坦
% * 在去掉尖点后通过 \(t=t^3/t^2\) 验证该映射为同构；分别说明“原点处非分歧失败”“原点处平坦失败”和“原点外为同构”三项结论
%
% ---
%

\subsection{尖点正规化的局部计算}
\label{b3:subsec:2304:01}

本节固定特征零域 \(k\)，考察平面尖点及其参数化
\[
A=k[t^2,t^3]\subseteq B=k[t].
\]
这是一个尤其适合比较局部性质的例子：参数环很简单，映射有限且双有理，但在尖点处既不非分歧，也不平坦。后三个小节分别计算微分模、特殊纤维和去掉尖点后的环，以区分这三项判断。

\begin{proposition}\label{b3:prop:cusp-normalization-setup}
设 \(k\) 为特征零域，\(A=k[t^2,t^3]\)、\(B=k[t]\)。则
\[
A\cong k[x,y]/(y^2-x^3),\qquad B=A+At,
\qquad\operatorname{Frac}(A)=\operatorname{Frac}(B)=k(t).
\]
并且 \(B\) 是 \(A\) 在这个分式域中的整闭包。尖点对应 \(\mathfrak m=(t^2,t^3)\subseteq A\)，其上方唯一的素理想为 \(\mathfrak n=(t)\subseteq B\)。
\end{proposition}
\begin{proof}
映射 \(k[x,y]\to B\)、\(x\mapsto t^2,y\mapsto t^3\) 的像为 \(A\)，核包含 \(y^2-x^3\)。模去此关系后，每个多项式都能写成 \(F(x)+yG(x)\)。代入后，\(F(t^2)\) 只含偶次幂，\(t^3G(t^2)\) 只含至少三次的奇次幂，不能相互抵消；所以核恰为上述主理想。

\(t\) 满足首一方程 \(T^2-t^2=0\)，且所有次数至少二的单项式都在 \(A\) 中，因此 \(B=A+At\) 为有限整扩张。又有 \(t=t^3/t^2\in\operatorname{Frac}(A)\)，故分式域相同。\(B=k[t]\) 为整闭整环；若分式域中元素整于 \(A\)，它也满足系数在 \(B\) 中的同一个首一方程，故属于 \(B\)。这证明整闭包的断言。

若 \(\mathfrak q\subseteq B\) 位于 \(\mathfrak m\) 上方，则 \(t^2\in\mathfrak q\)，从而 \(t\in\mathfrak q\)，所以 \(\mathfrak q=(t)\)。反向收缩显然为 \(\mathfrak m\)。
\end{proof}

局部环 \(A_{\mathfrak m}\) 维数为一，而 \(\mathfrak m/\mathfrak m^2\) 以 \(t^2,t^3\) 的像为基，维数为二，故不正则。相反，\(B_{\mathfrak n}=k[t]_{(t)}\) 是离散赋值环，因而正则。正规化改善了目标坐标环本身的奇点，但这种改善并不意味着原来的环同态光滑。

\subsection{光滑性、正则性与可分性}
\label{b3:subsec:2304:02}

正则性属于局部环本身：比较 \(\dim R\) 与 \(\dim_{R/\mathfrak m}\mathfrak m/\mathfrak m^2\) 就可作出判断。光滑性属于一个带有底环的同态。同一个域看作自身上的代数总是光滑，换成一个较小而不完美的基域以后却可能不光滑。

\begin{theorem}[完美基域上的比较；外部引用]\label{b3:thm:perfect-field-regular-smooth}
设 \(k\) 为完美域，\(B\) 为有限型交换 \(k\) 代数，\(\mathfrak q\in\operatorname{Spec}B\)。则 \(B_{\mathfrak q}\) 为正则局部环，当且仅当存在 \(b\notin\mathfrak q\)，使 \(k\to B_b\) 光滑。
\end{theorem}

这里引用松村英之的 Jacobian 判据及其比较说明\cite[Theorem 30.3 and Remark 2, pp. 233--235]{Matsumura1989CommutativeRingTheory}；相应的局部几何正则性刻画见\cite[Theorem 28.7, pp. 219--220]{Matsumura1989CommutativeRingTheory}。其中的局部判据结合有限组关系，把非零 Jacobian 子式固定在一个主开邻域中，便得到本定理的邻域表述。本章不另行重证该判据中的可分性部分。

\begin{example}\label{b3:exm:regular-not-smooth-imperfect}
设 \(k=\mathbb F_p(s)\)，\(L=k[u]/(u^p-s)\)。多项式 \(T^p-s\) 不可约，故 \(L\) 是域，自身作为零维局部环正则。但是 \(L\) 在 \(k\) 上不光滑。为直接看到提升失败，令
\[
C=k[T]/\bigl((T^p-s)^2\bigr),\qquad
J=(T^p-s)C.
\]
则 \(J^2=0\)、\(C/J\cong L\)。若恒等映射 \(L\to C/J\) 可以提升，\(u\) 的像必为 \(T+h\)、\(h\in J\)。然而特征 \(p\) 下
\((T+h)^p-s=T^p-s+h^p=T^p-s\ne0\)，因为 \(h^p=0\)。关系不成立，矛盾。

\ProseParagraphBreak
这个例子已在第14章出现；这里的新比较是 \(L\) 的正则性与 \(k\to L\) 的不光滑性可以同时成立。基变换后
\(L\otimes_kL\cong L[\varepsilon]/(\varepsilon^p)\)
也解释了差别：正则的原环在不可分基变换后产生幂零元。
\end{example}

上述完美性条件不能简单改写成“光滑点的剩余域必须可分于基域”。例如在 \(k=\mathbb F_p(s)\) 上，多项式代数 \(k[T]\) 光滑，但闭点 \((T^p-s)\) 的剩余域正是这个不可分扩张。光滑性研究整个代数映射的无穷小行为，并不要求每个闭点本身都是一个平展的 \(k\) 点。

\subsection{代数几何中的不同拓扑}
\label{b3:subsec:2304:03}

Zariski 邻域只使用环的局部化；平展映射则容许加入在给定点附近唯一可提升的代数解。为了使这些不同的局部研究方式各有相应的覆盖概念，代数几何使用比普通开集覆盖更一般的约定。这里只介绍仿射环所描述的三种覆盖，不展开层与上同调理论。

\begin{definition}\label{b3:def:affine-etale-fppf-covers}
设 \(A\) 为交换含幺环，\((A\to B_i)_{i\in I}\) 为一族交换 \(A\) 代数。若每个 \(\mathfrak p\in\operatorname{Spec}A\) 都是某个 \(B_i\) 中素理想的收缩，就称这一族在素谱上联合满射。
\begin{enumerate}
\item 形如 \(A\to A_{f_i}\) 且联合满射的族，是仿射的 Zariski 基本开覆盖；其条件等价于 \((f_i\mid i\in I)=A\)。
\item 若每个 \(A\to B_i\) 平展且这一族联合满射，称为一个\term[pingzhan-fugai]{平展覆盖}[étale covering]。
\item 若每个 \(A\to B_i\) 平坦、有限呈示且这一族联合满射，称为一个 \term[fppf-fugai]{fppf 覆盖}[fppf covering]。
\end{enumerate}
\end{definition}

字母 fppf 来自法语 \emph{fidèlement plat et de présentation finie}。对单个映射，平坦并在素谱上满射等价于忠实平坦，所以这个名称与第12章的代数条件相符。基本开覆盖是平展覆盖，平展覆盖也是 fppf 覆盖；但这些覆盖的每一项不必是原素谱的一个开子集。

\ProseParagraphBreak
例如平方映射 \(\mathbb C[s,s^{-1}]\to\mathbb C[t,t^{-1}]\) 是单个平展覆盖，因为它有限自由且素谱满射，却不是基本开嵌入。特征 \(p\) 下的映射 \(k[s]\to k[t]\)、\(s\mapsto t^p\) 则是 fppf 覆盖，而不是平展覆盖。因此在这些覆盖中说“局部存在解”，可能需要先扩张环；它不等同于在原来的 Zariski 开集中寻找解。

\ProseParagraphBreak
一般的平展拓扑与 fppf 拓扑把这样的覆盖作为基本资料，交叠由张量积描述。第12章的忠实平坦下降说明，扩张环以后得到的模及同构，只有满足交叠上的相容条件才来自原环。这也是为什么采用更丰富的覆盖后，仍须记录如何把局部对象拼回去。

\subsection{尖点正规化的微分模与非分歧性}
\label{b3:subsec:2304:04}

\begin{proposition}\label{b3:prop:cusp-relative-differentials}
设 \(k\) 为特征零域，\(A=k[t^2,t^3]\subseteq B=k[t]\)。则作为 \(B\) 模，
\[
\Omega_{B/A}\cong B/(t)\,dt.
\]
因此 \(A\to B\) 在 \(\mathfrak n=(t)\) 处不是非分歧映射。
\end{proposition}
\begin{proof}
对 \(k\to A\to B\) 的微分传递正合列为
\[
B\otimes_A\Omega_{A/k}\longrightarrow\Omega_{B/k}
\longrightarrow\Omega_{B/A}\longrightarrow0.
\]
\(\Omega_{B/k}=B\,dt\)，而 \(A\) 由 \(t^2,t^3\) 在 \(k\) 上生成，所以第一箭头的像由
\(d(t^2)=2t\,dt\)、\(d(t^3)=3t^2\,dt\) 生成。特征零使 \(2\) 可逆，故像为 \((t)\,dt\)，得到所述商模。在 \((t)\) 处局部化后，这个模仍同构于 \(k\,dt\)，不为零，所以该点不满足形式非分歧判据，也不满足非分歧条件。
\end{proof}

这一计算说明参数化在尖点处有一个不能被底环中的函数检测的一阶方向：对平方零参数 \(\varepsilon\)，代入 \(t=\varepsilon\) 后，\(t^2,t^3\) 都变成零，但 \(t\) 本身未必为零。因此底环只看到同一个点，上方却存在不同的无穷小变化。

\subsection{尖点纤维的长度与非平坦性}
\label{b3:subsec:2304:05}

\begin{proposition}\label{b3:prop:cusp-fiber-nonflat}
设 \(k\) 为特征零域，\(A=k[t^2,t^3]\subseteq B=k[t]\)，\(\mathfrak m=(t^2,t^3)\)、\(\mathfrak n=(t)\)。则特殊纤维为
\[
B\otimes_A A/\mathfrak m
\cong B/(t^2,t^3)B
\cong k[t]/(t^2),
\]
其长度为二；而一般纤维为 \(k(t)\)，维数为一。因此局部映射 \(A_{\mathfrak m}\to B_{\mathfrak n}\) 不平坦。
\end{proposition}
\begin{proof}
特殊纤维的同构直接由商与张量积相容得到；\(1,t\) 构成其 \(k\) 基，且唯一极大理想 \((t)\) 的平方为零，所以长度为二。令 \(K=\operatorname{Frac}(A)=k(t)\)。由于 \(t=t^3/t^2\)，有 \(B\otimes_AK=K\)，故一般秩为一。

注意 \(B\otimes_AA_{\mathfrak m}\) 是有限 \(A_{\mathfrak m}\) 代数，其极大理想只能位于 \(\mathfrak m\) 上方，而那里唯一的素理想是 \(\mathfrak n\)。因此这个有限代数是局部环，并等于 \(B_{\mathfrak n}\)：\(B\) 中 \(\mathfrak n\) 外元素在它里面全为单位。

若 \(B_{\mathfrak n}\) 在 \(A_{\mathfrak m}\) 上平坦，则它作为有限生成平坦模，由\cref{b3:lem:finite-flat-local-free}必须自由，设秩为 \(r\)。与 \(K\) 张量给出 \(r=1\)。再与剩余域 \(k\) 张量，便应得到一维 \(k\) 空间，却与刚算出的二维纤维矛盾。因此局部映射不平坦。
\end{proof}

这里比较纤维长度时，有限模在局部环上自由是必不可少的论证。单说“纤维大小发生变化”不足以证明一般映射不平坦；必须把这种变化与平坦性在当前有限条件下的具体后果联系起来。

\subsection{尖点外的同构与局部性质的比较}
\label{b3:subsec:2304:06}

\begin{proposition}\label{b3:prop:cusp-isomorphism-off-origin}
设 \(k\) 为特征零域，\(A=k[t^2,t^3]\subseteq B=k[t]\)。则
\[
A_{t^2}=B_{t^2}=k[t,t^{-1}].
\]
因此正规化映射在去掉尖点后为同构；尖点之外的局部映射都平展且光滑。
\end{proposition}
\begin{proof}
在 \(A_{t^2}\) 中，\(t=t^3/t^2\)，而 \(t\) 的逆元为 \(t/t^2\)，故 \(k[t,t^{-1}]\subseteq A_{t^2}\)。反向包含显然，\(B_{t^2}\) 也等于这个环。又若 \(A\) 的素理想包含 \(t^2\)，则由 \((t^3)^2=(t^2)^3\) 也包含 \(t^3\)，从而等于极大理想 \(\mathfrak m\)。所以 \(D(t^2)\) 正是尖点的补集。环同构保持全部平方零提升且作为模自由秩一，故局部映射平展和光滑。
\end{proof}

现在可以分别读出三个结论。在原点，\(\Omega_{B/A}\) 不为零，所以非分歧性失败；特殊纤维长度二而一般秩一，所以平坦性失败；在原点外，显式公式 \(t=t^3/t^2\) 给出同构。这些结论来自三个不同的计算，不能以“正规化去掉了奇点”一句话相互替代。

\ProseParagraphBreak
另一方面，\(B_{(t)}\) 本身正则，\(k\to B\) 也光滑，但 \(A_{\mathfrak m}\to B_{(t)}\) 不光滑，因为光滑必平坦。这个例子再次说明：讨论光滑时必须写清底环；讨论非分歧时必须保留无穷小信息；讨论平坦时则必须考察模的关系是否在张量之后仍被正确保留。


\begin{chaptersummary}
Hensel 提升首先是一个误差修正过程：单根条件使一元导数可逆，互素条件使控制因式修正的线性映射可逆，完备性再使各阶近似收敛。Hensel 化用泛性质只加入保证这些代数提升所需的元素；即使它与原环有相同的各阶商，也未必等于完备化。

\ProseParagraphBreak
平方零提升把关系的误差变成线性问题。形式光滑使关系模嵌入自由微分模并分裂；有限呈示进一步把这个分裂转化为邻域上的标准方程组。一般平坦性的证明还采用了明确列出的 Noether 局部定理。加上微分模消失以后，平展可以等价地描述为唯一提升，或局部标准平展模型。

\ProseParagraphBreak
尖点正规化把这些区别落实为计算：相对微分模为 \(k\,dt\)，特殊纤维为长度二的 \(k[t]/(t^2)\)，原点外则由 \(t=t^3/t^2\) 得到同构。正则的参数环、双有理性以及点集上的行为，都不能代替对给定环同态的光滑、非分歧和平坦条件逐项检验。
\end{chaptersummary}

\BookExercises

\begin{optionalexercise}\label{b3:ex:23-five-adic-root}
在 \(\mathbb Z_5\) 中求 \(T^2+1\) 同余于 \(2\pmod5\) 的根模 \(125\) 的值，并解释所得剩余类为何能继续唯一提升。
\end{optionalexercise}
\begin{solution}
先写 \(a=2+5c\)，模 \(25\) 得 \(5+20c\equiv0\)，故 \(c\equiv1\pmod5\)，即 \(a\equiv7\pmod{25}\)。再写 \(a=7+25d\)，模 \(125\) 得 \(50+350d\equiv0\)，即 \(2+4d\equiv0\pmod5\)，所以 \(d\equiv2\pmod5\)。因此所求值为 \(57\pmod{125}\)，确有 \(57^2+1=3250=26\cdot125\)。导数 \(2a\) 始终模 \(5\) 非零，单根提升定理保证存在唯一的 \(5\) 进根具有这个初始剩余类。
\end{solution}

\begin{optionalexercise}\label{b3:ex:23-series-root-factors}
在 \(\mathbb Q[[t]]\) 中求 \(X^2-X-t\) 同余于 \(0\pmod t\) 的根到 \(t^4\) 为止，并写出完整多项式的两个首一因子。解释为何这些因子不都属于 \(\mathbb Q[t]_{(t)}[X]\)。
\end{optionalexercise}
\begin{solution}
设根为 \(a=c_1t+c_2t^2+c_3t^3+c_4t^4+\cdots\)。比较 \(a^2-a-t=0\) 的各阶系数，依次得到
\(c_1=-1\)、\(c_2=c_1^2=1\)、\(c_3=2c_1c_2=-2\)、\(c_4=2c_1c_3+c_2^2=5\)。另一个根为 \(1-a\)，所以准确的分解是 \((X-a)(X-(1-a))\)，其中 \(a\) 表示完整的 Hensel 根。若 \(a\in\mathbb Q(t)\)，则判别式 \(1+4t\) 为有理函数的平方，与它的简单零点矛盾。因此两个一次因子的系数不能都来自原局部环。
\end{solution}

\begin{optionalexercise}\label{b3:ex:23-repeated-root-henselian}
设 \(k\) 为域，\(R=k[\varepsilon]/(\varepsilon^2)\)。证明 \(R\) 为 Hensel 环，而 \(T^2\) 模极大理想的重根 \(0\) 有多个提升。由此说明 Hensel 性没有取消单根条件。
\end{optionalexercise}
\begin{solution}
\(R\) 为局部环，极大理想 \((\varepsilon)\) 的平方为零，所以其进拓扑的逆极限从第二阶起稳定，\(R\) 完备分离，因而 Hensel。对每个 \(c\in k\)，元素 \(c\varepsilon\) 都同余于零且平方为零，至少 \(0\) 与 \(\varepsilon\) 不同。导数 \(2T\) 在约化根处为零，故唯一单根提升定理不适用。
\end{solution}

\begin{optionalexercise}\label{b3:ex:23-cubic-hensel-factor}
在 \(\mathbb Z_7[T]\) 中，将 \(T^3-T+7\) 分解为一个约化为 \(T\) 的首一次因子和一个约化为 \(T^2-1\) 的首二次因子，并求前一因子的常数项模 \(49\)。
\end{optionalexercise}
\begin{solution}
\(T\) 与 \(T^2-1\) 模 \(7\) 互素，所以唯一存在这样的分解。设同余于零的根为 \(a\)，则
\[
T^3-T+7=(T-a)\bigl(T^2+aT+a^2-1\bigr),
\]
因为余项为 \(a^3-a+7=0\)。写 \(a=7c\)，模 \(49\) 的根方程为 \(-7c+7\equiv0\)，所以 \(a\equiv7\pmod{49}\)，首一次因子的常数项为 \(-7\pmod{49}\)。二次因子的两个剩余根还可分别提升，但本题的因式提升不需要先选择它们。
\end{solution}

\begin{optionalexercise}\label{b3:ex:23-henselization-universal}
仅用 Hensel 化的泛性质证明：若局部环 \(R\) 已经 Hensel，则 \(R\to R^{\mathrm h}\) 是同构；两个 Hensel 化之间的相容同构唯一。再解释为什么相同的有限阶商不能证明一个环已经完备。
\end{optionalexercise}
\begin{solution}
若 \(R\) Hensel，恒等映射给出唯一局部同态 \(v:R^{\mathrm h}\to R\)，使 \(v\iota=1_R\)。两个从 \(R^{\mathrm h}\) 到自身并延拓 \(\iota\) 的局部同态 \(1\) 与 \(\iota v\) 由唯一性相同，故互逆。对两个 Hensel 化，分别使用泛性质得到相容映射，再用同样的复合论证即得唯一同构。至于有限阶商，正文中的 \(\mathbb Q[t]_{(t)}\)、它的 Hensel 化及 \(\mathbb Q[[t]]\) 有相同的各阶商，但前两者没有包含全部相容级数；完备性要求环到这些商的逆极限为同构，单独比较每个商没有证明满射性。
\end{solution}

\begin{optionalexercise}\label{b3:ex:23-node-lifting-obstruction}
设 \(k\) 为域，\(B=k[x,y]/(xy)\)。用一个平方零提升问题证明 \(k\to B\) 不形式光滑。
\end{optionalexercise}
\begin{solution}
取 \(C=k[\varepsilon]/(\varepsilon^3)\)、\(J=(\varepsilon^2)\)。定义 \(u:B\to C/J\) 将 \(x,y\) 都送到 \(\varepsilon\)，这是同态，因为乘积在 \(C/J\) 中为零。任意提升都必须把二者分别送到 \(\varepsilon+a\varepsilon^2\)、\(\varepsilon+b\varepsilon^2\)，而乘积恒为非零的 \(\varepsilon^2\)。关系不能成立，因此没有提升。失败发生在同时满足关系的要求，而不是变量像自身不能选取提升。
\end{solution}

\begin{optionalexercise}\label{b3:ex:23-free-differentials-not-smooth}
设 \(k\) 为特征二域，\(B=k[x]/(x^2)\)。证明 \(\Omega_{B/k}\) 为秩一自由模，但 \(B\) 不形式光滑。指出关系模分裂判据中的哪一步失败。
\end{optionalexercise}
\begin{solution}
微分关系为 \(d(x^2)=2x\,dx=0\)，所以 \(\Omega_{B/k}=B\,dx\)。取 \(P=k[X]\)、\(I=(X^2)\)，则 \(I/I^2\) 为以 \(X^2\) 的像为基的秩一自由 \(B\) 模，但到 \(B\,dX\) 的映射为零，没有左逆。也可取 \(C=k[T]/(T^4)\)、\(J=(T^2)\)，恒等映射 \(B\to C/J\) 的任何候选提升把 \(x\) 送到 \(T+h\)、\(h\in J\)，其平方为 \(T^2\ne0\)，故没有提升。这表明微分模投射只是形式光滑的一个后果，不能独自代替关系模的分裂。
\end{solution}

\begin{optionalexercise}\label{b3:ex:23-etale-stability}
证明有限呈示平展映射在复合和任意基变换下保持平展。证明中应说明有限呈示条件为何也得到保持。
\end{optionalexercise}
\begin{solution}
先使用平展等价于有限呈示且形式平展。形式平展在复合与基变换下保持已在正文证明。基变换把有限个生成元与有限条关系的系数映到新底环，仍给出有限呈示。对复合 \(A\to B\to C\)，取 \(B\) 的有限生成元和关系，再添上 \(C\) 在 \(B\) 上的有限生成元；后者有限条关系所含的 \(B\) 系数都能写成前一组生成元的多项式。把这些关系一并写在 \(A\) 上，得到有限呈示。因此复合与基变换仍满足平展的等价条件。
\end{solution}

\begin{optionalexercise}\label{b3:ex:23-hyperbola-standard-model}
设 \(A\) 为交换含幺环，\(B=A[x,y]/(xy-1)\)。给出标准光滑模型，计算 \(\Omega_{B/A}\)，并判断当 \(A\ne0\) 时映射是否平展。
\end{optionalexercise}
\begin{solution}
环中 \(x,y\) 互逆，所以 Jacobian 行 \((y,x)\) 中的 \(x\) 已为单位。以变量 \(y\) 对应的列为子式，得到相对维数一的标准光滑模型。微分关系为 \(y\,dx+x\,dy=0\)，故 \(dy=-x^{-2}dx\)，而 \(\Omega_{B/A}=B\,dx\)。若 \(A\ne0\)，则 \(B\cong A[x,x^{-1}]\ne0\)，这个微分模不为零，所以映射光滑且平坦，却不平展。
\end{solution}

\begin{optionalexercise}\label{b3:ex:23-gaussian-integers-etale}
对 \(\mathbb Z\to\mathbb Z[i]\) 计算相对微分模，证明该映射平坦，并准确给出其平展的点。
\end{optionalexercise}
\begin{solution}
\(\mathbb Z[i]\cong\mathbb Z[T]/(T^2+1)\)，所以作为整数模以 \(1,i\) 为基，因而平坦。相对微分模为
\(\mathbb Z[i]/(2i)\,di=\mathbb Z[i]/(2)\,di\)，因为 \(i\) 可逆。故在 \(D(2)\) 上导数可逆，映射平展；在包含 \(2\) 的素理想处该模非零，不平展。\(2=-i(1+i)^2\)，唯一这样的素理想是 \((1+i)\)。这显示有限自由不自动非分歧，整数 \(2\) 必须保留在判断中。
\end{solution}

\begin{optionalexercise}\label{b3:ex:23-closed-point-unramified}
设 \(A=k[t]\)、\(B=A/(t)=k\)。证明 \(\Omega_{B/A}=0\)，但 \(B\) 不是平坦 \(A\) 模，并直接给出形式光滑性失败的平方零问题。
\end{optionalexercise}
\begin{solution}
对商代数，所有元素都来自 \(A\)，所以任意 \(A\) 导子为零，\(\Omega_{B/A}=0\)。\(A\) 上的单射乘法 \(A\xrightarrow{t}A\) 与 \(B\) 张量后成为零映射 \(B\to B\)，不是单射，故不平坦。取 \(C=A/(t^2)\)、\(J=(t)C\)，则 \(C/J=B\)。若恒等映射 \(B\to C/J\) 有 \(A\) 代数提升，\(t\) 在 \(B\) 中为零会迫使其在 \(C\) 中的像为零，但它实际非零。这说明唯一性条件可以成立而存在性条件失败。
\end{solution}

\begin{optionalexercise}\label{b3:ex:23-product-diagonal}
令 \(C=A\times A\)，其中 \(A\to C\) 为对角同态。用 \(e_1=(1,0)\)、\(e_2=(0,1)\) 写出正文对角分裂中的元素 \(z\)，并显式分裂任意 \(C\) 模的映射 \(C\otimes_AM\to M\)。
\end{optionalexercise}
\begin{solution}
取 \(z=e_1\otimes e_1+e_2\otimes e_2\)。乘法把它送到 \(e_1+e_2=1\)，而对每个 \(c=(a,b)\)，有 \((c\otimes1)z=(1\otimes c)z\)。所以截面为
\(m\mapsto e_1\otimes e_1m+e_2\otimes e_2m\)。它为 \(C\) 线性且复合乘法为恒等。对应的分解是 \(M=e_1M\oplus e_2M\)；若 \(M\) 在 \(A\) 上平坦，两个直和项均平坦，这也直接证明 \(M\) 在 \(A\times A\) 上平坦。
\end{solution}

\begin{optionalexercise}\label{b3:ex:23-finite-flat-quotient}
设 \(C\) 为交换含幺环，\(J\subseteq C\) 为有限生成理想。证明 \(C/J\) 在 \(C\) 上平坦，当且仅当 \(J=(e)\) 对某个幂等元 \(e\in C\) 成立。解释相应闭子谱为何也为开子谱。
\end{optionalexercise}
\begin{solution}
若商平坦，用它张量 \(J\hookrightarrow C\) 得到单射 \(J/J^2\to C/J\)，但映射为零，故 \(J=J^2\)。正文的行列式技巧给出 \(J=(e)\)、\(e^2=e\)。反之，\(C=eC\oplus(1-e)C\)，且 \(C/J\cong(1-e)C\) 为自由模 \(C\) 的直和项，故投射并平坦。对每个素理想，\(e(1-e)=0\) 且 \(e+(1-e)=1\)，所以恰有一个幂等因子属于该素理想。因此 \(V(e)=D(1-e)\)，既闭又开。
\end{solution}

\begin{optionalexercise}\label{b3:ex:23-frobenius-fiber-length}
设 \(k\) 为特征 \(p>0\) 的代数闭域，\(k[s]\to k[t]\) 由 \(s\mapsto t^p\) 给出。计算点 \(s=a\) 上的纤维及其长度，说明幂零纤维为何并不与平坦性矛盾。
\end{optionalexercise}
\begin{solution}
唯一取 \(b\in k\) 使 \(b^p=a\)，则纤维为
\(k[t]/(t^p-a)=k[t]/((t-b)^p)\cong k[\varepsilon]/(\varepsilon^p)\)，长度为 \(p\)。原扩张以 \(1,t,\ldots,t^{p-1}\) 为自由基，秩正是 \(p\)，所以特殊纤维长度与自由秩一致，完全符合平坦性。它有幂零元说明非分歧性失败；平坦性本身不要求纤维约化。
\end{solution}

\begin{optionalexercise}\label{b3:ex:23-cusp-other-characteristics}
撤去尖点计算中的特征零假设，对任意域 \(k\) 计算
\(\Omega_{k[t]/k[t^2,t^3]}\)。分别给出特征二、特征三及其他特征的答案，并判断原点外是否仍为同构。
\end{optionalexercise}
\begin{solution}
传递正合列仍给出
\(\Omega_{B/A}=B/(2t,3t^2)\,dt\)。特征二时，\(2=0\)、\(3=1\)，所以商为 \(B/(t^2)\,dt\)；特征三时，\(2\) 可逆，商为 \(B/(t)\,dt\)；其他特征下也因 \(2\) 可逆得到 \(B/(t)\,dt\)。在特征二中原点处的微分模长度为二，其余情形长度为一，均不为零。原点外的公式 \(t=t^3/t^2\) 不涉及特征，故同构结论始终成立。
\end{solution}

\begin{optionalexercise}\label{b3:ex:23-cusp-module-presentation}
设 \(A=k[t^2,t^3]\)、\(B=k[t]\)，记 \(x=t^2,y=t^3\)。证明作为 \(A\) 模，\(B\) 有展示
\[
A^2\xrightarrow{\begin{psmallmatrix}-y&-x^2\\x&y\end{psmallmatrix}}
A^2\xrightarrow{(1\quad t)}B\longrightarrow0.
\]
利用这个展示重算尖点纤维的维数。
\end{optionalexercise}
\begin{solution}
后一箭头满射，因为 \(B=A+At\)。两列确实给出关系：\(-y+xt=0\)、\(-x^2+yt=0\)。反之，若 \(a+bt=0\)、\(a,b\in A\)，则 \(bt\in A\)。\(A\) 的多项式恰好没有一次项，所以 \(b\) 的常数项必须为零，即 \(b\in(x,y)\)。写 \(b=xc+yd\)，便有 \(a=-yc-x^2d\)，因此 \((a,b)\) 是所列两列的线性组合，证明正合。模去 \((x,y)\) 后矩阵为零，张量积的右正合性给出纤维为 \(k^2\)，与 \(k[t]/(t^2)\) 的基 \(1,t\) 一致。
\end{solution}

\begin{optionalexercise}\label{b3:ex:23-inseparable-closed-point}
设 \(k=\mathbb F_p(s)\)，\(B=k[T]\)，\(\mathfrak q=(T^p-s)\)。证明 \(k\to B\) 光滑、\(B_{\mathfrak q}\) 正则，但剩余域 \(B/\mathfrak q\) 不可分于 \(k\)。解释这为何不与正文比较定理冲突。
\end{optionalexercise}
\begin{solution}
多项式环有任意变量像的提升，所以 \(k\to B\) 光滑。\(T^p-s\) 不可约，\(B_{\mathfrak q}\) 是主理想整环在非零素理想处的局部化，为 DVR，故正则。剩余域 \(L=k[u]\)、\(u^p=s\)，其极小多项式导数为零，故扩张纯不可分且次数为 \(p\)。正文没有断言光滑的每个点都具有可分剩余域；本例的整个多项式代数仍有一个可自由提升的变量，而商掉 \(T^p-s\) 所得的 \(k\to L\) 已经是另一个映射，后者不光滑。
\end{solution}

\begin{optionalexercise}\label{b3:ex:23-square-root-cover}
设 \(A=\mathbb C[s,s^{-1}]\)，\(B=A[T]/(T^2-s)\)。证明 \(\operatorname{Spec}B\to\operatorname{Spec}A\) 是平展覆盖，但在任何非空 Zariski 开集上都没有与底环相容的截面。说明“平展局部存在平方根”在这个例子中的含义。
\end{optionalexercise}
\begin{solution}
\(B\) 有自由基 \(1,T\)，且 \(T\) 为单位，所以导数 \(2T\) 可逆，映射平展。自由秩二的模忠实平坦，因而素谱映射满射，构成平展覆盖。若在某个非空开集上有截面，限制到一般点便给出 \(\mathbb C(s)\) 中的元素 \(r\)，满足 \(r^2=s\)。但平方的零点阶数为偶数，而 \(s\) 在零点有一阶零点，矛盾。所谓平展局部存在，是说先经过这个平展覆盖，即把底环换成 \(B\)，平方根 \(T\) 就已经在新环中；并不是说原来的某个非空 Zariski 开集已经含有这样的有理函数。
\end{solution}
